\section{Related work}\label{sec:relatedworks}
{\color{blue}
    Credit risk prediction has been extensively studied in the literature using a wide range of pattern recognition and machine learning algorithms \cite{gasmi2025features, bai2025improved, dougan2025credit, aruleba2025improved, yao2025novel, liang2025application, lu2025efficient, yu2025hybrid, nguyen2025comparative, zou2025variable, zeng2025nexus, zhuang2024design, talaat2024toward, mosa2024ccfd, hartomo2025novel, ali2025explainable, zandi2025attention, pavitha2024explainable, jacobs2024benchmarking, mienye2023deep, bakhtiari2023credit, gao2023cate, sun2025enhancing, zhou2025enhancing, baisholan2025fraudx, chen2024interpretable, de2022explainable, ariza2020explainability, moscato2021benchmark, nallakaruppan2024credit, nobel2024unmasking, do2024explainable, zhao2024resampling, chang2024credit}. However, only some of these works have addressed the explainability of the prediction, a critical aspect in financial applications where predictions must be transparent and justifiable \cite{sun2025enhancing, zhou2025enhancing, ali2025explainable, hartomo2025novel, baisholan2025fraudx, chen2024interpretable, de2022explainable, ariza2020explainability, moscato2021benchmark, nallakaruppan2024credit, nobel2024unmasking, do2024explainable, chang2024credit}. In the following, we review key works in the literature on credit risk prediction, with a particular focus on recent approaches that either incorporate explainability or achieve high predictive performance. Accordingly, the review is divided into two parts: works that focus solely on predictive accuracy without offering explanations, and works that integrate explainability into their predictors.\\ %In the following, we review the main works published in the literature on the most recent ones that include some form of explainability. Accordingly, the review is divided into two parts: works focused on achieving high predictive performance without providing explanations, and works that integrate explainability into predictions. \\%Additionally, we consider works that achieve strong predictive performance but lack explainability, to highlight the existing trade-off between accuracy and explainability, and to motivate the need for explainable alternatives.\\
    
    % Primero empezaré con aquellos que no explican pero tiene buenos resultados
    In \cite{aruleba2025improved}, a stacked ensemble predictor that combines predictors such as Random Forests, Logistic Regression, and Convolutional Neural Networks, together with SMOTE-ENN \cite{batista2004study} to balance imbalanced datasets. While \cite{oliveira2025advancing} proposes the use of artificial neural networks for credit risk prediction, these outperform other predictors such as Logistic Regression, Support Vector Machines, Decision Trees, Random Forests, and gradient boosting.\\ 

    In \cite{lin2025deep}, the authors integrate a neural network with a Special Forces Algorithm (SFA) for hyperparameter optimization in credit risk prediction, achieving superior performance compared to XGBoost, LightGBM, and LSTM. In contrast, \cite{sundaravadivel2025optimizing} employed a Random Forest classifier in combination with SMOTE \cite{chawla2002smote} to address class imbalance in the dataset for credit risk prediction.\\

    % Después seguire con los trabajos que explain of prediction
    % zhou2025enhancing: ya quedo al implementación
    % do2024explainable: ya quedo al implementación
    %In \cite{zhou2025enhancing}, the authors proposed the use of XGBoost for credit risk prediction and SMOTE for a balanced dataset. For explanation of the predictions, the authors integrated the SHAP, which assigns a contribution to each feature and shows how each feature influences the prediction. %The predictor achieved high performance in terms of precision, recall, and F1-score. 
    %Similarly, \cite{nobel2024unmasking} the authors use XGBoost for credit risk prediction, SMOTE for a balanced dataset, and SHAP to explain the prediction. 
    In \cite{zhou2025enhancing} and \cite{nobel2024unmasking}, the authors use XGBoost for credit risk prediction and SMOTE to address class imbalance. Both studies use SHAP to generate explanations that highlight the contribution of each feature to the prediction. On the other hand, \cite{do2024explainable}, \cite{hernes2023credit}, and \cite{chang2024credit} address credit risk prediction using XGBoost and integrate the SHAP to explain the predictions. In particular, \cite{chang2024credit} proposes a variant of SHAP designed to produce explanations in the presence of correlation between features. % ya quedo al implementación
    %In \cite{baisholan2025fraudx}, the authors proposed to combine Random Forest and XGBoost for credit risk prediction. To explain the predictions, they incorporated the SHAP to highlight the contribution of each input feature to the prediction. Similarly, \cite{sun2025enhancing} proposed combining Random Forest and XGBoost for credit risk prediction. To explain prediction, the authors applied the SHAP to show the contribution of each feature to the prediction.\\
    In \cite{baisholan2025fraudx} and \cite{sun2025enhancing}, the authors propose combining Random Forest and XGBoost for credit risk prediction, evaluating their approaches on different datasets. Both studies used SHAP to explain predictions to highlight the contribution of each feature. \\%These works are compared with predictors as Random Forest, Logistic Regression, Decision Tree, Multi Layer Perceptron, LightGBM, and Gradient Boosting.\\
    
    In \cite{nallakaruppan2024explainable}, the authors propose using Random Forest with explainable AI techniques for credit risk prediction. To generate explanations for predictions, LIME, SHAP, and Partial Dependence Plots (PDPs) were used, which help identify how feature values contribute to the prediction. Similarly, \cite{nguyen2025comparative} presents a comparative study of boosting algorithms, as Gradient Boosting, AdaBoost, XGBoost, and CatBoost, for credit risk prediction. To address the class imbalance problem often present in credit datasets, they employed SMOTETomek \cite{batista2003balancing}. The authors analyze the importance of features and employ SHAP to explain the predictions.\\

    %In \cite{hernes2023credit}, the authors propose XGBoost for credit risk prediction. To explain predictions, the authors used SHAP, ceteris paribus plots, and feature importance. Experimental results indicate that XGBoost outperforms Logistic Regression in most performance metrics, except for sensitivity, where Logistic Regression shows a higher rate of correct predictions for credit risk prediction. \\    

    In \cite{zandi2025attention}, the authors present a multilayer dynamic graph neural network for credit risk prediction. The prediction for each instance is based on its features and the evolution of their relationships over time, using SHAP to explain the prediction through feature importance. In \cite{talaat2024toward}, the authors propose CreditNetXAI, which combines deep neural networks with XAI techniques for credit risk prediction, using feature selection methods to identify the most relevant features for prediction. To explain predictions, they use SHAP to assign an importance to each feature and indicate its contribution to the prediction.\\ 

    In \cite{zheng2024micro}, the authors employ Decision Tree-based LightGBM with an optimized gradient approach to credit risk prediction. To explain the prediction, they use SHAP to visualize the importance of features in the prediction. Similarly in \cite{de2022explainable} used a LightGBM for credit risk prediction and SHAP to explain the prediction. In \cite{han2024nate}, the authors propose NATE, a non-parametric approach to credit risk prediction in imbalanced datasets that integrates gradient boosting as a predictor. NATE utilised SMOTE to balance the dataset and employed SHAP to explain the prediction, highlighting the contribution of each feature to the prediction.\\

    All of these studies evaluated the predictive performance of their predictors using standard classification metrics such as accuracy, precision, recall, F1-score, and, in some cases, the type II error. These metrics were applied to evaluate the performance of each predictor in credit risk prediction, especially in the context of imbalanced datasets. As discussed in previous works, explanations in credit risk prediction have largely focused on assigning importance scores to individual features using SHAP and LIME. While these approaches improve explanation, they often fail to capture the underlying structure of the data, particularly the recurring combinations of feature values that characterize different types of customers. This limitation is critical, as financial institutions increasingly demand not only high predictive accuracy but also clear, instance-specific justifications to comply with regulatory requirements and enhance customer communication.\\

    To address this challenge, our approach introduces a pattern-based explanation framework. By mining frequent feature-value patterns from the training data and associating them with clusters of similar instances, we provide structured explanations that reflect common behaviors observed in historical cases. These patterns offer a more intuitive and human-understandable rationale for predictions.